\section{Conclusion}
\label{sec:conclusion}

We introduce \silobench{} to evaluate distributed coordination in multi-agent LLM systems across 1,620 experiments. 
The results are unambiguous: current LLMs cannot reliably escape their information silos through coordination alone.

The \textit{Communication-Reasoning Gap} identifies the precise fault line: agents are competent communicators but poor distributed reasoners. 
They spontaneously form task-appropriate topologies and exchange information actively, yet consistently fail to integrate what they have gathered---a dissociation made concrete by the PCS--SR divergence and the RCC analysis showing that coordination overhead eliminates parallelization gains entirely at high complexity. 
Most strikingly, spontaneous leader emergence actively hurts performance on complex tasks, revealing that self-organized centralization creates bottlenecks rather than resolving them.

Closing this gap will require mechanisms for information sufficiency detection, explicit consensus protocols, and adaptive coordination strategies. 
\silobench{} provides the evaluation infrastructure to track progress along these directions.
